\section{Cauchy 积分公式}\label{ch10:section-cauchy}

复变函数论研究的核心定理之一便是解析函数的 Cauchy 积分公式. 对于这个定理的经典形式而言, 作为预备知识我们不仅需要 Jordan 曲线定理, 并且还需要上一节中介绍的环绕数定理. Cauchy 积分公式有一个另类表达方式, 它不需要用到这些结论. 尽管这个表达方式使人感到有些别扭, 但却经常出现在有关的教材之中.

由于我们已经掌握了 Jordan 曲线定理, 所以我们给自己规定的任务是从那个另类表达方式出发来推导经典形式的 Cauchy 积分公式.

我们先更为正式地介绍``环绕数''这个概念.

\begin{definition}[winding number]\label{ch10:def-winding-number}
设 $f$ 是 $\mathbb{R}^{2}$ 中的一条回路, 点 $a$ 不是 $f$ 的像点. 令
\[
g(s)=\frac{[f(s)-a]}{\|f(s)-a\|},
\]
那么 $g$ 是 $S^{1}$ 中的一条回路. 令 $p\colon\mathbb{R}\to S^{1}$ 是标准覆叠映射, $\widetilde{g}$ 为 $g$ 到 $\mathbb{R}$ 的提升. 由于 $g$ 是一条回路, 所以差 $\widetilde{g}(1)-\widetilde{g}(0)$ 是整数. 这个整数就称为 $f$ 关于 $a$ 的环绕数 (winding number of $f$ with respect to $a$), 记为 $n(f,a)$.
\end{definition}

注意 $n(f,a)$ 并不依赖于 $g$ 的提升的选取. 因为如果 $\widetilde{g}$ 是 $g$ 的一个提升, 那么提升的唯一性蕴涵着 $g$ 的任意其他提升可以表示为 $\widetilde{g}(s)+m$, 其中 $m$ 是某个整数.

\begin{definition}[free homotopy]\label{ch10:def-free-homotopy}
设 $F\colon I\times I\to X$ 是一个连续映射, 对任意 $t$ 有 $F(0,t)=F(1,t)$. 那么对于每一个 $t$, 映射 $f_{t}(s)=F(s,t)$ 是 $X$ 中的一条回路. 映射 $f$ 称为回路 $f_{0}$ 和 $f_{1}$ 之间的一个自由同伦 (free homotopy). 自由同伦是回路之间的一个同伦, 在同伦的过程中回路的基点允许移动.
\end{definition}

\begin{lemma}\label{ch10:lemma-winding-properties}
设 $f$ 是 $\mathbb{R}^{2}-a$ 中的一条回路.

\begin{enumerate}[(a)]
\item 若 $\overline{f}$ 是 $f$ 的逆, 则 $n(\overline{f},a)=-n(f,a)$.
\item 若 $f$ 通过 $\mathbb{R}^{2}-a$ 中的回路自由同伦于 $f'$, 则 $n(f,a)=n(f',a)$.
\item 若 $a$ 和 $b$ 属于 $\mathbb{R}^{2}-f(I)$ 的同一分支, 则 $n(f,a)=n(f,b)$.
\end{enumerate}
\end{lemma}

\begin{proof}
(a) 为了计算 $n(\overline{f},a)$, 整个定义中我们用 $1-s$ 代替 $s$. 这恰好改变了 $\widetilde{g}(1)-\widetilde{g}(0)$ 的符号.

(b) 设 $F$ 是 $f$ 和 $f'$ 之间的自由同伦. 定义 $G\colon I\times I\to S^{1}$ 为
\[
G(s,t)=\frac{[F(s,t)-a]}{\|F(s,t)-a\|}.
\]
令 $\widetilde{G}$ 是 $G$ 到 $\mathbb{R}$ 的提升. 则对于每一个 $t$, $\widetilde{G}(1,t)-\widetilde{G}(0,t)$ 是一个整数. 又根据连续性, 它是一个常数.

(c) 设 $\alpha$ 是 $\mathbb{R}^{2}-f(I)$ 中从 $a$ 到 $b$ 的道路. 根据定义, $n(f,a)=n(f-a,0)$. 因为 $f(s)-\alpha(t)$ 是 $\mathbb{R}^{2}-0$ 中 $f-a$ 和 $f-b$ 之间的自由同伦, 由此我们的结论成立.
\end{proof}

\begin{definition}[simple loop]\label{ch10:def-simple-loop}
设 $f$ 是 $X$ 中的一条回路. 我们称 $f$ 是简单回路 (simple loop), 如果只在 $s=s'$ 或 $s,s'$ 中一个为 $0$, 另一个为 $1$ 的情况下才有 $f(s)=f(s')$. 如果 $f$ 是一条简单回路, 那么它的像集便是 $X$ 中的一条简单闭曲线.
\end{definition}

\begin{theorem}\label{ch10:thm-winding-simple}
设 $f$ 是 $\mathbb{R}^{2}$ 中的一条简单回路. 若 $a$ 属于 $\mathbb{R}^{2}-f(I)$ 的一个无界分支, 则 $n(f,a)=0$. 若 $a$ 属于一个有界分支, 则 $n(f,a)=\pm 1$.
\end{theorem}

\begin{proof}
由于 $n(f,a)=n(f-a,0)$, 我们可以只讨论 $a=0$ 这种情形. 进而, 可以假定 $f$ 的基点在正 $x$ 轴上. 因为我们可以旋转 $\mathbb{R}^{2}-0$ 使得 $f$ 的基点到达这种位置, 旋转过程中 $f$ 的改变过程是一个自由同伦的过程, 所以对于定理的结论没有影响.

设 $f$ 是 $X=\mathbb{R}^{2}-0$ 中的一条简单回路, 其基点 $x_{0}$ 位于正 $x$ 轴上. 设 $C$ 为简单闭曲线 $f(I)$. 以下证明, 当 $0$ 在 $\mathbb{R}^{2}-C$ 的有界分支中时, $[f]$ 生成 $\pi_{1}(X,x_{0})$. 当 $0$ 在 $\mathbb{R}^{2}-C$ 的无界分支中时, $[f]$ 是平凡的.

通过标准的商映射 $p\colon I\to S^{1}$, 映射 $f$ 诱导出一个同胚 $h\colon S^{1}\to C$. 元素 $[p]$ 生成 $S^{1}$ 的基本群, 所以 $h_{*}[p]$ 生成 $C$ 的基本群. 如果 $0$ 在 $\mathbb{R}^{2}-C$ 的有界分支中, 定理 65.2 告诉我们 $j_{*}h_{*}[p]=[f]$ 生成 $\mathbb{R}^{2}-0$ 的基本群, 其中 $j\colon C\to\mathbb{R}^{2}-0$ 是内射. 另一方面, 如果 $0$ 在 $\mathbb{R}^{2}-C$ 的无界分支中, 则根据零伦引理, $j\circ h$ 是零伦的, 所以 $[f]$ 是平凡的.

我们来证明: 若 $[f]$ 生成 $\pi_{1}(X,x_{0})$, 则 $n(f,0)=\pm 1$. 若 $[f]$ 是平凡的, 则 $n(f,0)=0$. 由于从 $\mathbb{R}^{2}-0$ 到 $S^{1}$ 上的收缩 $x\mapsto x/\|x\|$ 诱导基本群之间的同构, 回路 $g(s)=f(s)/\|f(s)\|$ 在第一种情形下代表 $\pi_{1}(S^{1},b_{0})$ 的一个生成元, 而在第二种情形下则代表单位元. 考察圆周基本群定理证明中的同构 $\phi\colon\pi_{1}(S^{1},b_{0})\to\mathbb{Z}$, 我们便会发现当将 $g$ 提升为 $\mathbb{R}$ 中的一条以 $0$ 为起点的道路 $\widetilde{g}$ 时, 在第一种情形下, 道路 $\widetilde{g}$ 将在 $\pm 1$ 处终结. 而在第二种情形下, 将在 $0$ 处终结.
\end{proof}

\begin{definition}[counterclockwise / clockwise]\label{ch10:def-orientation}
设 $f$ 是 $\mathbb{R}^{2}$ 中的一条简单回路. 称 $f$ 是逆时针 (counterclockwise) 回路, 如果对于 $\mathbb{R}^{2}-f(I)$ 的有界分支中的某一个点 (因此对于每一个点) $a$ 有 $n(f,a)=+1$. 称 $f$ 为顺时针 (clockwise) 回路, 如果 $n(f,a)=-1$. 因此标准回路 $p(s)=(\cos 2\pi s,\sin 2\pi s)$ 是逆时针回路.
\end{definition}

\subsubsection*{在复变函数中的应用}

我们现在来讨论环绕数与复积分之间的关联.

\begin{lemma}\label{ch10:lemma-winding-integral}
设 $f$ 是复平面上的一条分段可微的回路, $a$ 是一个不在 $f$ 的像中的点. 则
\[
n(f,a)=\frac{1}{2\pi\mathrm{i}}\int_{f}\frac{\mathrm{d}z}{z-a}.
\]
这个等式常用来定义 $f$ 的环绕数.
\end{lemma}

\begin{proof}
这个引理的证明是一个简单的计算练习. 设 $p\colon\mathbb{R}\to S^{1}$ 是标准覆叠映射. 令 $r(s)=\|f(s)-a\|$, $g(s)=[f(s)-a]/r(s)$. 设 $\widetilde{g}$ 是 $g$ 到 $\mathbb{R}$ 上的提升, $\theta(s)=2\pi\widetilde{g}(s)$. 则 $f(s)-a=r(s)\exp(\mathrm{i}\theta(s))$, 因此
\[
\begin{aligned}
\int_{f}\frac{\mathrm{d}z}{z-a}
&=\int_{0}^{1}\frac{(r'\mathrm{e}^{\mathrm{i}\theta}+\mathrm{i}r\theta'\mathrm{e}^{\mathrm{i}\theta})/r\mathrm{e}^{\mathrm{i}\theta}\,\mathrm{d}s}{}\\
&=[\log r(s)+\mathrm{i}\theta(s)]_{0}^{1}\\
&=\mathrm{i}[\theta(1)-\theta(0)]\\
&=2\pi\mathrm{i}[\widetilde{g}(1)-\widetilde{g}(0)].
\end{aligned}
\]
\end{proof}

\begin{theorem}[Cauchy integral formula---classical version]\label{ch10:thm-cauchy-classical}
设 $C$ 是复平面上的一条分段可微的简单闭曲线, $B$ 是 $\mathbb{R}^{2}-C$ 的一个有界分支. 如果 $F(z)$ 在包含 $B$ 和 $C$ 的开集 $\Omega$ 上是解析的, 那么对于 $B$ 中的每一个点 $a$ 都有
\[
F(a)=\pm\frac{1}{2\pi\mathrm{i}}\int_{C}\frac{F(z)}{z-a}\,\mathrm{d}z.
\]
其中, 若 $C$ 是逆时针定向的, 上式中的符号取``$+$'', 反之则取``$-$''.
\end{theorem}

\begin{proof}
我们从 Ahlfors $[\mathrm{A}]$ 中证明的形式出发来推导这个公式. 在 Ahlfors $[\mathrm{A}]$ 中陈述的定理是:

设 $F$ 在区域 $\Omega$ 上是解析的, $f$ 在 $\Omega$ 上是一条分段可微的回路. 假设对不在 $\Omega$ 中的每个点 $b$, 有 $n(f,b)=0$. 若 $a\in\Omega$ 并且 $a$ 不是 $f$ 的像点, 则
\[
n(f,a)\cdot F(a)=\frac{1}{2\pi\mathrm{i}}\int_{f}\frac{F(z)}{z-a}\,\mathrm{d}z.
\]

现在我们将这个结果用到简单闭曲线 $C$ 的一个分段可微参数化 $f$ 上. 当 $b$ 不在 $\Omega$ 中时, $n(f,b)=0$ 成立, 因为每一个这样的 $b$ 都属于 $\mathbb{R}^{2}-C$ 的无界分支. 进而, 当 $a\in B$ 时, 根据定理 66.2, $n(f,a)=\pm 1$, 其符号依赖于 $C$ 的定向. 定理证毕.
\end{proof}

注意, 如果没有关于 Jordan 曲线定理的知识, 我们甚至不能将 Cauchy 积分定理的经典形式陈述清楚. 为了给出证明, 我们还要用到关于简单闭曲线的环绕数的知识. 有趣的是, 后一结果可以用一个完全不同的方法来证明 (至少, 对于可微的情形是如此). 这个方法就是应用在分析中证明的 Green 定理的一般形式. 我们在习题 2 中给出了一个证明概要.

\begin{exercise}\label{ch10:ex-cauchy-1}
设 $f$ 是 $\mathbb{R}^{2}-a$ 中的一条回路, $g(s)=[f(s)-a]/\|f(s)-a\|$. 映射 $g$ 通过标准商映射 $p\colon I\to S^{1}$ 诱导出一个连续映射 $h\colon S^{1}\to S^{1}$. 证明 $n(f,a)$ 等于 $h$ 的度, 其定义见第 58 节中的习题 9.
\end{exercise}

\begin{exercise}\label{ch10:ex-cauchy-2}
这个练习要求熟悉流形分析.

\textbf{定理.} 设 $C$ 是 $\mathbb{R}^{2}$ 中的一条简单闭曲线, 并且是 $\mathbb{R}^{2}$ 中的光滑子流形. 设 $f\colon I\to C$ 是一条简单回路, 它光滑地参数化 $C$. 若 $0$ 是 $\mathbb{R}^{2}-C$ 的有界分支的一个点, 那么 $n(f,0)=\pm 1$.

证明: 设 $U$ 是 $\mathbb{R}^{2}-C$ 的有界分支, $B$ 是 $U$ 中的一个以 $0$ 为圆心的 $\varepsilon$-闭球. 设 $S=\operatorname{Bd}B$, $M$ 是 $U-B$ 的闭包.

\begin{enumerate}[(a)]
\item 证明 $M$ 是一个以 $C\cup S$ 为边界的 $2$-维光滑流形.
\item 应用 Green 定理证明 $\int_{C}\mathrm{d}z/z=\pm\int_{S}\mathrm{d}z/z$, 其中符号依赖于 $S$ 和 $C$ 的定向. [提示: 令 $P=-y/(x^{2}+y^{2})$, $Q=x/(x^{2}+y^{2})$.]
\item 证明第二个积分等于 $\pm 2\pi\mathrm{i}$.
\end{enumerate}
\end{exercise}